\section*{Hoofdstuk \ref{ch:b2:structures} --- Verzamelingen en structuren}

\begin{solution}{exo:b2:structures:1}
\emph{Eindige deelverzamelingen van $\N$:} \omterm{def:b2:structures:countable}{aftelbaar} --- de
verzameling van de deelverzamelingen van $\intint{0}{n}$ is eindig,
en de eindige deelverzamelingen vormen de \omterm{def:b2:structures:countable}{aftelbare} vereniging
daarvan over $n$
(\cref{prop:b2:structures:countablestable}~(3)); oneindig, want zij
bevat alle eenpuntsverzamelingen.

\emph{Alle deelverzamelingen van $\N$:} niet \omterm{def:b2:structures:countable}{aftelbaar}, wegens de
stelling van Cantor (\cref{thm:b2:structures:cantor}~(1) met $E =
\N$).

\emph{$\R \setminus \Q$:} niet \omterm{def:b2:structures:countable}{aftelbaar} --- anders zou $\R = \Q
\cup (\R\setminus\Q)$ een vereniging van twee \omterm{def:b2:structures:countable}{aftelbare
verzamelingen} zijn, in strijd met
\cref{thm:b2:structures:cantor}~(2).

\emph{Veeltermen over $\Q$:} \omterm{def:b2:structures:countable}{aftelbaar} --- de veeltermen van graad
$\leq n$ injecteren in $\Q^{n+1}$ (eindige producten van \omterm{def:b2:structures:countable}{aftelbare
verzamelingen}), en neem de vereniging over $n$.

\emph{Binaire rijen die vanaf zeker moment nul zijn:} \omterm{def:b2:structures:countable}{aftelbaar} ---
zij staan in bijectie met de eindige deelverzamelingen van $\N$
(hun drager).
\end{solution}

\begin{solution}{exo:b2:structures:2}
Reken element voor element, met de rechterfactor eerst. Nu stuurt
$\sigma\tau$ de elementen $1 \mapsto \sigma(1) = 4$, $\;2 \mapsto
\sigma(3) = 5$, $\;3 \mapsto \sigma(7) = 7$, $\;4 \mapsto
\sigma(4) = 2$, $\;5 \mapsto \sigma(5) = 3$, $\;6 \mapsto
\sigma(6) = 1$, $\;7 \mapsto \sigma(2) = 6$:
\[
\sigma\tau = (1\,4\,2\,5\,3\,7\,6),
\]
een $7$-cykel. Evenzo stuurt $\tau\sigma$ de elementen $1 \mapsto
\tau(4) = 4$, $\;2 \mapsto \tau(6) = 6$, $\;3 \mapsto \tau(5) =
5$, $\;4 \mapsto \tau(2) = 3$, $\;5 \mapsto \tau(3) = 7$, $\;6
\mapsto \tau(1) = 1$, $\;7 \mapsto \tau(7) = 2$:
\[
\tau\sigma = (1\,4\,3\,5\,7\,2\,6),
\]
eveneens een $7$-cykel (zoals te verwachten was: $\sigma\tau$ en
$\tau\sigma$ zijn geconjugeerd en hebben dus hetzelfde cykeltype).

\omterm{def:b2:structures:generated}{Orden} en signaturen: $\sigma$ heeft cykeltype $(4,2)$: \omterm{def:b2:structures:generated}{orde}
$\operatorname{lcm}(4,2) = 4$, signatuur $(-1)^3(-1)^1 = +1$;
$\tau$ is een $3$-cykel: \omterm{def:b2:structures:generated}{orde} $3$, signatuur $+1$; beide producten
zijn $7$-cykels: \omterm{def:b2:structures:generated}{orde} $7$, signatuur $(-1)^6 = +1$.

$\sigma^{2026}$: uit $2026 = 4 \times 506 + 2$ volgt
$\sigma^{2026} = \sigma^2 = (1\,2)(4\,6)$ (kwadrateer de
$4$-cykel; de \omterm{def:b2:structures:sn}{transpositie} kwadrateert weg).
\end{solution}

\begin{solution}{exo:b2:structures:3}
$\{0,1\}^{\N} \to \mathcal{P}(\N)$: een rij gaat naar haar drager
--- een bijectie (indicatorfuncties), geen stelling nodig.

$\{0,1\}^{\N} \to \intcc{0}{1}$: de afbeelding in grondtal $3$,
$(a_n) \mapsto \sum 2a_n 3^{-n-1}$, is injectief (twee
verschillende rijen verschillen voor het eerst op rang $N$; de
staarten kunnen een sprong van $2\cdot 3^{-N-1}$ niet compenseren,
want $\sum_{n > N} 2\cdot 3^{-n-1} = 3^{-N-1} < 2\cdot3^{-N-1}$).

$\intcc{0}{1} \to \{0,1\}^{\N}$: de binaire ontwikkeling, waarbij
we (zeg) de ontwikkeling kiezen die niet op louter $1$'en eindigt:
injectief.

Cantor--Bernstein (\cref{thm:b2:structures:cantorbernstein})
toegepast op de laatste twee injecties geeft dat $\intcc{0}{1}$ en
$\{0,1\}^{\N}$ \omterm{def:b2:structures:countable}{gelijkmachtig} zijn, en dus alle drie de
verzamelingen.
\end{solution}

\begin{solution}{exo:b2:structures:4}
Zij $c = ab = ba$ en $d = \operatorname{ord}(c)$. Om te beginnen is
$c^{mn} = a^{mn} b^{mn} = e$ (dankzij de commutatie mag de macht
worden gesplitst), dus $d \mid mn$. Omgekeerd geeft $c^d = e$ dat
$a^d = b^{-d}$; dit element ligt in $\langle a\rangle \cap \langle
b\rangle$, een ondergroep waarvan de \omterm{def:b2:structures:generated}{orde} zowel $m$ als $n$ deelt
(Lagrange in elk van beide \omterm{def:b2:structures:generated}{cyclische groepen}) en die dus triviaal
is: $a^d = b^d = e$, zodat $m \mid d$ en $n \mid d$, en wegens de
onderlinge ondeelbaarheid $mn \mid d$. Bijgevolg is $d = mn$.

In $\mathfrak{S}_3$: neem $a = (1\,2)$ (\omterm{def:b2:structures:generated}{orde} $2$) en $b =
(1\,2\,3)$ (\omterm{def:b2:structures:generated}{orde} $3$), met onderling ondeelbare \omterm{def:b2:structures:generated}{orden}, die niet
commuteren: $ab = (2\,3)$ heeft \omterm{def:b2:structures:generated}{orde} $2 \neq 6$ ---
$\mathfrak{S}_3$ heeft immers geen element van \omterm{def:b2:structures:generated}{orde} $6$. De
commutatie is dus onmisbaar.
\end{solution}

\begin{solution}{exo:b2:structures:5}
Koppel elke $x \in G$ aan $x^{-1}$. De paren $\{x, x^{-1}\}$ met
$x \neq x^{-1}$ tellen twee elementen en verdelen hun vereniging;
de overige elementen zijn precies die met $x = x^{-1}$, dat wil
zeggen $x^2 = e$. Omdat $\abs G$ even is en de tweetallige paren
een even aantal elementen bedekken, heeft de verzameling $\{x : x^2
= e\}$ een even aantal elementen; zij bevat $e$, dus bevat zij
minstens nog één ander element $x \neq e$ --- een element van \omterm{def:b2:structures:generated}{orde}
$2$.
\end{solution}

\begin{solution}{exo:b2:structures:6}
Elk element van $A_n$ is een product van een even aantal
\omterm{def:b2:structures:sn}{transposities} (\cref{thm:b2:structures:signature}: ontbind in
\omterm{def:b2:structures:sn}{transposities}; het aantal is even omdat de signatuur $+1$ is). Het
volstaat dus elk product van twee \omterm{def:b2:structures:sn}{transposities} met $3$-cykels te
schrijven:
\[
(a\,b)(a\,c) = (a\,c\,b),
\qquad
(a\,b)(c\,d) = (a\,c\,b)(a\,c\,d) \quad (\text{verschillende }
a,b,c,d),
\]
(na te gaan door te evalueren), en $(a\,b)(a\,b) = \mathrm{id}$.
De $3$-cykels brengen dus $A_n$ voort.
\end{solution}

\begin{solution}{exo:b2:structures:7}
\emph{$(\Q,+) \to (\Z,+)$:} alleen het nulmorfisme. Voor elke $x$
en elke $n \geq 1$ is $f(x) = n f\bigl(\frac xn\bigr)$ deelbaar
door $n$ in $\Z$; het enige gehele getal dat door elke $n$
deelbaar is, is $0$, dus $f(x) = 0$ voor alle $x$.

\emph{$(\Z/n\Z, +) \to (\Z/m\Z, +)$:} een morfisme ligt vast door
$c = f(\overline 1)$, dat aan $n c \equiv 0 \pmod m$ moet voldoen,
dat wil zeggen: $c$ is een veelvoud van $\frac{m}{\gcd(m,n)}$; er
zijn $\gcd(m,n)$ zulke klassen, en elke keuze definieert ook
werkelijk een morfisme (factoriseer $k \mapsto kc$ over $\Z/n\Z$
met de universele eigenschap).

\emph{$(\Q, +) \to (\Q_+^*, \times)$:} alleen het triviale. Is
$f(x) = y$, dan is $y = f(n \cdot \frac xn) = f(\frac xn)^n$ voor
elke $n$ een $n$-de macht in $\Q_+^*$. Maar een rationaal getal $y
\neq 1$ kan niet voor alle $n$ een $n$-de macht zijn: een zeker
priemgetal komt in $y$ voor met een exponent $v \neq 0$, en $n
\nmid v$ zodra $n > \abs v$ (exponenten van $n$-de machten zijn
veelvouden van $n$, wegens de eenduidige factorisatie). Bijgevolg
is $f \equiv 1$.
\end{solution}

\begin{solution}{exo:b2:structures:8}
Uit $360 = 2^3 \cdot 3^2 \cdot 5$ volgt
$\varphi(360) = 360\bigl(1 - \tfrac12\bigr)\bigl(1 -
\tfrac13\bigr)\bigl(1 - \tfrac15\bigr) = 360 \cdot \tfrac12 \cdot
\tfrac23 \cdot \tfrac45 = 96$.

Het stelsel: de moduli $8, 9, 5$ zijn paarsgewijs onderling
ondeelbaar, met product $360$. Uit $x \equiv 3 \pmod 8$ en $x
\equiv 5 \pmod 9$: $x = 3 + 8k$ met $3 + 8k \equiv 5 \pmod 9$, dus
$-k \equiv 2$, $k \equiv -2 \equiv 7 \pmod 9$: $x \equiv 3 + 56 =
59 \pmod{72}$. Vervolgens $59 + 72\ell \equiv 2 \pmod 5$: $4 +
2\ell \equiv 2$, $2\ell \equiv 3 \equiv 8$, $\ell \equiv 4 \pmod
5$: $x \equiv 59 + 288 = 347 \pmod{360}$.

De laatste twee cijfers van $3^{2026}$: modulo $4$ is $3^{2026} =
9^{1013} \equiv 1$. Modulo $25$: $\varphi(25) = 20$ en $2026 =
20\cdot101 + 6$, dus $3^{2026} \equiv 3^6 = 729 \equiv 4
\pmod{25}$. Los $x \equiv 1 \pmod 4$, $x \equiv 4 \pmod{25}$ op:
uit $x = 4 + 25k \equiv 1 \pmod 4$ volgt $k \equiv 1 \pmod 4$,
dus $x \equiv 29 \pmod{100}$. De laatste twee cijfers zijn $29$.
\end{solution}

\begin{solution}{exo:b2:structures:9}
Zij $A$ een eindig integriteitsgebied en $a \in A$ met $a \neq 0$.
De afbeelding $x \mapsto ax$ is injectief ($ax = ay \implies a(x -
y) = 0 \implies x = y$, want er zijn geen nuldelers); een
injectieve afbeelding van een eindige verzameling naar zichzelf is
surjectief (volume van bachelorjaar 1, de equivalentie met het
duivenhokprincipe). Dus is $1 = ab$ voor zekere $b$: elk element
ongelijk aan nul is inverteerbaar en $A$ is een lichaam.

Voor $\Z/n\Z$: is $n$ priem, dan is het een integriteitsgebied
($n \mid ab \implies n \mid a$ of $n \mid b$, het lemma van
Euclides), eindig, en dus een lichaam; is $n = rs$ samengesteld,
dan levert $\overline r\,\overline s = \overline 0$ nuldelers.
\end{solution}

\begin{solution}{exo:b2:structures:10}
Zij $m = \max\{\operatorname{ord}(x) : x \in G\}$, aangenomen in
$a$.

\emph{Bewering: de \omterm{def:b2:structures:generated}{orde} van elke $x \in G$ deelt $m$.} Stel dat een
zekere $x$ \omterm{def:b2:structures:generated}{orde} $q$ heeft met $q \nmid m$: dan deelt een zekere
priemmacht $p^k$ wel $q$ maar niet $m$. Schrijf $m = p^j m'$ met $p
\nmid m'$ en $j < k$. Het element $a^{p^j}$ heeft \omterm{def:b2:structures:generated}{orde} $m'$; het
element $x^{q/p^k}$ heeft \omterm{def:b2:structures:generated}{orde} $p^k$; deze \omterm{def:b2:structures:generated}{orden} zijn onderling
ondeelbaar en de twee elementen commuteren ($G \subseteq K^*$ is
abels), dus heeft hun product volgens \cref{exo:b2:structures:4}
\omterm{def:b2:structures:generated}{orde} $p^k m' > p^j m' = m$: in strijd met de maximaliteit.

Alle $x \in G$ voldoen dus aan $x^m = 1$: de veelterm $X^m - 1$
heeft minstens $\abs G$ nulpunten in het lichaam $K$, waaruit
$\abs G \leq m$ volgt (een veelterm van graad $m$ ongelijk aan nul
heeft hoogstens $m$ nulpunten, volume van bachelorjaar 1). Maar
$m = \operatorname{ord}(a) \leq \abs G$ volgens Lagrange. Dus is $m
= \abs G$ en is $\langle a \rangle$, met $m = \abs G$ elementen,
heel $G$: \omterm{def:b2:structures:generated}{cyclisch}.

Voor $K = \Z/p\Z$: $(\Z/p\Z)^*$ is een eindige ondergroep van
$K^*$ en dus \omterm{def:b2:structures:generated}{cyclisch} (van \omterm{def:b2:structures:generated}{orde} $p - 1$).
\end{solution}

\begin{solution}{exo:b2:structures:11}
\emph{Niet \omterm{def:b2:structures:generated}{cyclisch}:} de ondergroep $\langle \frac pq\rangle$
bestaat uit de gehele veelvouden van $\frac pq$, en die hebben
alle een noemer die $q$ deelt (in vereenvoudigde vorm); zij mist
dus $\frac{1}{2q}$. Geen enkele voortbrenger kan de onbegrensde
noemers van $\Q$ bereiken.

\emph{Niet eindig \omterm{def:b2:structures:generated}{voortgebracht}:} de ondergroep \omterm{def:b2:structures:generated}{voortgebracht} door
$\frac{p_1}{q_1}, \dots, \frac{p_k}{q_k}$ bestaat uit rationale
getallen waarvan de noemer $Q = q_1 \cdots q_k$ deelt (gehele
combinaties hebben een noemer die $Q$ deelt): zij mist
$\frac{1}{2Q}$.

\emph{Eindig \omterm{def:b2:structures:generated}{voortgebrachte} ondergroepen zijn \omterm{def:b2:structures:generated}{cyclisch}:} met $Q$
als hierboven ligt de ondergroep $H = \langle \frac{p_1}{q_1},
\dots, \frac{p_k}{q_k}\rangle$ in $\frac{1}{Q}\Z$. De afbeelding
$x \mapsto Qx$ is een isomorfisme van $\frac1Q\Z$ op $\Z$ dat $H$
naar een ondergroep van $\Z$ brengt, en die is $n\Z$ voor zekere
$n$ (volume van bachelorjaar 1): dus is $H = \frac{n}{Q}\Z$
\omterm{def:b2:structures:generated}{cyclisch}, \omterm{def:b2:structures:generated}{voortgebracht} door $\frac nQ$.
\end{solution}

\begin{solution}{exo:b2:structures:12}
\emph{Een \omterm{def:b2:structures:countable}{aftelbare} deelverzameling.} Zij $E$ oneindig. Construeer
$a_0, a_1, a_2, \dots$ inductief: $E$ is niet leeg, kies $a_0 \in
E$; zijn $a_0, \dots, a_n$ gekozen, dan is $E \setminus \{a_0,
\dots, a_n\}$ niet leeg ($E$ is niet eindig), kies daarin
$a_{n+1}$. De $a_n$ zijn per constructie paarsgewijs verschillend,
dus is $A = \{a_n : n \in \N\}$ een \omterm{def:b2:structures:countable}{aftelbare} deelverzameling van
$E$.

\emph{Oneindig $\implies$ \omterm{def:b2:structures:countable}{gelijkmachtig} met een echte
deelverzameling.} Definieer $f \colon E \to E \setminus \{a_0\}$
door $f(a_n) = a_{n+1}$ en $f(x) = x$ voor $x \notin A$. Zij is
injectief (beide stukken zijn injectief met disjuncte beelden) en
surjectief op $E \setminus \{a_0\}$: elke $a_{n+1}$ wordt bereikt
en elke $x \notin A$ ook. Dus is $E$ \omterm{def:b2:structures:countable}{gelijkmachtig} met de echte
deelverzameling $E \setminus \{a_0\}$.

\emph{Omkering.} Is $E$ eindig en $g \colon E \to F$ een bijectie
op $F \subseteq E$ met $F \neq E$, dan is $g$ een injectie van $E$
in zichzelf die niet surjectief is, in strijd met het
duivenhokprincipe (volume van bachelorjaar 1: een injectieve
afbeelding van een eindige verzameling naar zichzelf is
bijectief). Een verzameling die \omterm{def:b2:structures:countable}{gelijkmachtig} is met een echte
deelverzameling, is dus oneindig.
\end{solution}

\begin{solution}{pb:b2:structures:1}
\textbf{1.} Een stand kent aan elk van de $16$ vakjes precies één
van de $16$ inhouden toe (de plaatjes $1$--$15$ of de lege plek $b
= 16$), elk precies één keer: dat is juist een bijectie
$\intint1{16} \to \intint1{16}$, een element van
$\mathfrak{S}_{16}$; er zijn er $16! = 20\,922\,789\,888\,000$.
Een toegestane zet schuift één plaatje dat aan de lege plek grenst,
dus is het aantal zetten gelijk aan het aantal buren van het vakje
van de lege plek: $2$ voor de vier hoekvakjes, $3$ voor de acht
randvakjes en $4$ voor de vier inwendige vakjes.

\textbf{2.} Na het schuiven draagt vakje $p$ de vroegere inhoud van
$c$ en draagt vakje $c$ de lege plek; alle andere vakjes blijven
onaangeroerd: $\sigma'(p) = \sigma(c)$, $\sigma'(c) = \sigma(p) =
16$ en $\sigma' = \sigma$ elders. Dat is precies $\sigma' = \sigma
\circ (p\ c)$. Omdat $\varepsilon$ een morfisme is en
$\varepsilon\bigl((p\ c)\bigr) = -1$, volgt
$\varepsilon(\sigma') = -\varepsilon(\sigma)$.

\textbf{3.} Buurvakjes verschillen precies in één van de twee
coördinaten één stap, dus verandert $i + j$ van pariteit: $\chi$
neemt op buurvakjes tegengestelde waarden aan. Een zet brengt de
lege plek van $p$ naar het aangrenzende $c$ en klapt daarmee
$\chi(\text{vakje van de lege plek})$ om. Langs een gesloten
wandeling van de lege plek wordt $\chi$ per zet één keer omgeklapt
en keert zij terug naar haar beginwaarde: het aantal zetten is
even.

\textbf{4.} Volgens de vragen 2 en 3 klapt één zet beide factoren
van $I(\sigma) = \varepsilon(\sigma)\chi(\sigma^{-1}(16))$ om; hun
product verandert dus niet. Voor de opgeloste stand:
$\varepsilon(\mathrm{id}) = +1$ en de lege plek ligt in vakje $16$,
rij $4$, kolom $4$: $\chi(16) = (-1)^{8} = +1$, dus
$I(\mathrm{id}) = +1$.

\textbf{5.} $\sigma_L$ is de \omterm{def:b2:structures:sn}{transpositie} $(14\ 15)$ van vakjes:
$\varepsilon(\sigma_L) = -1$; haar lege plek is thuis, $\chi(16) =
+1$: $I(\sigma_L) = -1 \neq +1 = I(\mathrm{id})$. Omdat $I$ door
elke zet bewaard blijft, verbindt geen enkele reeks zetten
$\sigma_L$ met $\mathrm{id}$. De prijs was structureel veilig.

\textbf{6.} Leg een vakje $p$ vast en twee andere vakjes $c \neq d$
verschillend van $p$, en zet $\tau_0 = (c\ d)$. Op de verzameling
van de standen met de lege plek in $p$ is $\sigma \mapsto \sigma
\circ \tau_0$ een involutie (zij bewaart $\sigma(p) = 16$ omdat
$\tau_0$ het vakje $p$ vasthoudt) die $\varepsilon$ omklapt en dus
ook $I$: zij koppelt de standen met $I = +1$ bijectief aan die met
$I = -1$. Elk van de $16$ posities van de lege plek levert dus
$15!/2$ standen met $I = +1$, en
\[
\abs{\{I = +1\}} = 16 \cdot \frac{15!}{2} = \frac{16!}{2}.
\]

\textbf{7.} De zet die het plaatje van $c$ naar $p$ schuift, wordt
ongedaan gemaakt door datzelfde plaatje (nu in $p$) terug naar $c$
te schuiven: twee keer met $(p\ c)$ samenstellen geeft de
identiteit. Daaruit volgen reflexiviteit (de lege reeks),
symmetrie (draai de reeks om en maak elke zet ongedaan) en
transitiviteit (plak de reeksen aan elkaar): een
equivalentierelatie. Elke $\sigma \in R$ voldoet volgens vraag 4
aan $I(\sigma) = I(\mathrm{id}) = +1$, dus $R \subseteq \{I =
+1\}$; en $\sigma_L \notin R$ levert een tweede klasse.

\textbf{8.} Is $\sigma(16) = 16$, dan permuteert $\sigma$ de
vakjes $1, \dots, 15$; noem $\rho$ die beperking. Een vast punt
toevoegen verandert noch het cykeltype noch de signatuur (ontbind
$\rho$ in \omterm{def:b2:structures:sn}{transposities}; hetzelfde product werkt in
$\mathfrak{S}_{16}$), dus $\varepsilon(\sigma) =
\varepsilon(\rho)$, en met $\chi(16) = +1$ krijgen we $I(\sigma) =
\varepsilon(\rho)$. Elke stand kan naar een stand met de lege plek
thuis worden gebracht: het rooster is samenhangend, dus wandel de
lege plek langs een pad van buurvakjes naar vakje $16$ (elke stap
is een toegestane zet). Neem nu aan dat elke even $\rho \in
\mathfrak{S}_{15}$ door een programma wordt gerealiseerd. Zij
$\sigma$ gegeven met $I(\sigma) = +1$: wandel de lege plek naar
huis en bereik zo $\widetilde\sigma$ (equivalent met $\sigma$),
met $I(\widetilde\sigma) = +1$, dat wil zeggen: haar beperking
$\rho$ is even; het programma dat $\rho$ realiseert brengt
$\widetilde\sigma$ naar $\widetilde\sigma \circ \rho^{-1} =
\mathrm{id}$ (zie vraag 9). Wegens de transitiviteit is $\sigma
\in R$, waaruit $\{I = +1\} \subseteq R$ en dus de gelijkheid.

\textbf{9.} \emph{Eén zet:} de inhoud van $c$ belandt in $p$ en de
lege plek in $c$: het effect is $\pi = (p\ c)$, en inderdaad is
$\sigma' = \sigma \circ (p\ c) = \sigma \circ \pi^{-1}$.
\emph{Inductie:} heeft een reeks effect $\pi_1$ en brengt zij
$\sigma$ naar $\sigma \circ \pi_1^{-1}$, dan geeft een
daaropvolgende zet met effect $\pi_2 = (p'\ c')$ het resultaat
$(\sigma \circ \pi_1^{-1}) \circ \pi_2^{-1} = \sigma \circ
(\pi_2\pi_1)^{-1}$, en de inhouden verplaatsen zich volgens $\pi_2
\circ \pi_1$ (eerst $\pi_1$, dan $\pi_2$). Effecten stellen zich
dus samen, en een programma dat vanuit $\sigma$ wordt uitgevoerd
eindigt in $\sigma \circ \pi^{-1}$. \emph{Ondergroep:} het lege
programma heeft effect $\mathrm{id}$; aaneenschakelen geeft
producten; een programma omdraaien (vraag 7) geeft inversen. Het
effect van een programma houdt vakje $16$ vast (de lege plek
begint en eindigt thuis), dus $H \leq \mathfrak{S}_{15}$.
\emph{Pariteit:} een programma van $k$ zetten heeft $k$ even
(vraag 3), en $\varepsilon(\sigma \circ \pi^{-1}) =
(-1)^k\varepsilon(\sigma)$ dwingt $\varepsilon(\pi) = +1$: $H
\subseteq A_{15}$.

\textbf{10.} Volg de vier schuifbewegingen vanuit de lege plek in
$16$: de zet $16 \to 12$ brengt de inhoud van $12$ naar $16$; de
zet $12 \to 11$ brengt de inhoud van $11$ naar $12$; de zet $11
\to 15$ brengt de inhoud van $15$ naar $11$; de zet $15 \to 16$
brengt de inhoud die in $16$ geparkeerd stond (oorspronkelijk die
van $12$) naar $15$. Netto: $11 \mapsto 12$, $12 \mapsto 15$, $15
\mapsto 11$, lege plek thuis: het effect is $(11\ 12\ 15)$. De
omgekeerde rondgang maakt dit ongedaan: effect $(11\ 12\
15)^{-1} = (11\ 15\ 12)$. Beide zijn effecten van programma's en
liggen dus in $H$.

\textbf{11.} Aangrenzendheid van de opeenvolgende vakjes: binnen
elk genoemd paar verschillen de vakjes $1$ binnen dezelfde rij
($16{-}15$, $15{-}14$, $14{-}13$; $1{-}2$, $2{-}3$, $3{-}4$;
$8{-}7$, $7{-}6$; $10{-}11$, $11{-}12$) of $4$ binnen dezelfde
kolom ($13{-}9$, $9{-}5$, $5{-}1$; $4{-}8$; $6{-}10$; $12{-}16$):
een gesloten wandeling door alle $16$ vakjes, van lengte $16$.
Effect: net als in vraag 10, met de bezochte vakjes $c_0 = 16, c_1
= 15, \dots, c_{15} = 12$, gaat de inhoud van $c_i$ naar $c_{i-1}$
voor $i = 2, \dots, 15$, en wordt de inhoud van $c_1$, na de
eerste zet geparkeerd in $16$, door de laatste zet naar $c_{15}$
gebracht. Het effect stuurt dus $15 \mapsto 12$, en $14 \mapsto
15$, $13 \mapsto 14$, $9 \mapsto 13$, $5 \mapsto 9$, $1 \mapsto
5$, $2 \mapsto 1$, $3 \mapsto 2$, $4 \mapsto 3$, $8 \mapsto 4$, $7
\mapsto 8$, $6 \mapsto 7$, $10 \mapsto 6$, $11 \mapsto 10$, $12
\mapsto 11$: precies de $15$-cykel $\zeta$. Haar cyclische
volgorde begint met $x_0 = 15$, $x_1 = 12$, $x_2 = 11$, en $(x_0\
x_1\ x_2) = (15\ 12\ 11)$ stuurt $15 \mapsto 12 \mapsto 11 \mapsto
15$ --- en dat is precies $(11\ 15\ 12)$, de omgekeerde
elementaire rondgang.

\textbf{12.} Zij $\gamma = (a\ b\ c)$ en $x \in \intint1n$. Is $x
= g(a)$, dan $g\gamma g^{-1}(x) = g(\gamma(a)) = g(b)$; evenzo
gaat $g(b) \mapsto g(c)$ en $g(c) \mapsto g(a)$. Is $x \notin
\{g(a), g(b), g(c)\}$, dan ligt $g^{-1}(x) \notin \{a,b,c\}$ vast
onder $\gamma$, dus ligt ook $x$ vast. Bijgevolg is $g\gamma
g^{-1} = (g(a)\ g(b)\ g(c))$. En voor $g, h \in H$ is $ghg^{-1}
\in H$ volgens de ondergroepaxioma's.

\textbf{13.} Er geldt $\zeta \in H$ (vraag 11) en $s_0 = (x_0\ x_1\
x_2) \in H$ (vragen 10--11). Omdat $\zeta(x_i) = x_{i+1}$ (indices
modulo $15$), geeft vraag 12
\[
\zeta^{t}\,s_0\,\zeta^{-t}
= \bigl(\zeta^t(x_0)\ \zeta^t(x_1)\ \zeta^t(x_2)\bigr)
= (x_t\ x_{t+1}\ x_{t+2}) = s_t \in H
\qquad (t = 0, 1, \dots, 14).
\]

\textbf{14.} Op inverteren na mogen we aannemen dat $s = (a\ b\ c)$
en $t = (b\ c\ d)$ (een $3$-cykel op $\{a,b,c\}$ is $(a\ b\ c)$ of
haar inverse; evenzo op $\{b,c,d\}$; een voortbrenger door zijn
inverse vervangen laat $\langle s, t\rangle$ ongemoeid). Met $t$
eerst toegepast is dan
\[
st \colon a \mapsto b,\quad b \mapsto a,\quad c \mapsto d,\quad
d \mapsto c, \qquad\text{dus}\quad st = (a\ b)(c\ d),
\]
een dubbele \omterm{def:b2:structures:sn}{transpositie}. De ondergroep $G = \langle s, t\rangle$
bestaat uit even permutaties van de vier letters, dus $G \leq A_4$
en $\abs G \mid 12$; zij bevat een element van \omterm{def:b2:structures:generated}{orde} $3$ en een van
\omterm{def:b2:structures:generated}{orde} $2$, dus $6 \mid \abs G$ (Lagrange,
\cref{thm:b2:structures:lagrange}, toegepast op de twee cyclische
ondergroepen). Had $A_4$ een ondergroep $K$ van \omterm{def:b2:structures:generated}{orde} $6$, dan zou
die index $2$ hebben, en dan zou $g^2 \in K$ gelden voor elke $g
\in A_4$: voor $g \in K$ is dat duidelijk; voor $g \notin K$ zijn
$K$ en $gK$ de enige nevenklassen, dus is $g^2K$ gelijk aan $K$ of
aan $gK$, en $g^2K = gK$ zou $g \in K$ afdwingen. Elk kwadraat ligt
dus in $K$. Maar elke $3$-cykel $\gamma$ is een kwadraat, $\gamma =
(\gamma^2)^2$, en $A_4$ bevat acht $3$-cykels: $8 > 6$,
tegenspraak. Bijgevolg is $\abs G = 12$, dus $G = A_4$.

\textbf{15.} Breid $u \mapsto a$, $v \mapsto b$ uit tot een
bijectie $g_0$ van $X$ (stuur de overige $k - 2$ letters bijectief
waarheen dan ook op het complement van $\{a, b\}$). Is $g_0$
oneven, kies dan twee verschillende letters $s_1, t_1 \in X
\setminus \{u, v\}$ (mogelijk want $k \geq 4$) en vervang $g_0$
door $g_0 \circ (s_1\ t_1)$, dat even is en nog steeds $u \mapsto
a$, $v \mapsto b$ stuurt. Breid buiten $X$ uit met de identiteit:
een even permutatie $g \in G$ (het is een even permutatie van
$X$). Vraag 12 geeft dan
\[
g\,(u\ v\ w)\,g^{-1} = (g(u)\ g(v)\ g(w)) = (a\ b\ w) \in G,
\]
met gebruik van $g(w) = w$.

\textbf{16.} Elke $3$-cykel van $X \cup \{w\}$ ligt in $G$: die
met drager binnen $X$ zijn even permutaties van $X$; een met
drager $\{a, b, w\}$ is $(a\ b\ w)$ of $(b\ a\ w)$, en beide
worden door vraag 15 geleverd. Volgens \cref{exo:b2:structures:6}
brengen de $3$-cykels van de $(k+1)$-elementige verzameling $X
\cup \{w\}$ haar alternerende groep voort, dus bevat $G$ elke even
permutatie van $X \cup \{w\}$. \emph{Aaneenschakelen:} zij $G =
\langle s_0, \dots, s_{12}\rangle$. Lemma A toegepast op $s_0 =
(x_0\ x_1\ x_2)$ en $s_1 = (x_1\ x_2\ x_3)$ (de dragers delen
$\{x_1, x_2\}$) geeft alle even permutaties van $X_4 = \{x_0, x_1,
x_2, x_3\}$. Bevat $G$ alle even permutaties van $X_m = \{x_0,
\dots, x_{m-1}\}$ ($4 \leq m \leq 14$), dan heeft $s_{m-2} =
(x_{m-2}\ x_{m-1}\ x_m)$ de letters $u = x_{m-2}, v = x_{m-1} \in
X_m$ en de nieuwe letter $w = x_m$: Lemma B en het eerste deel
geven alle even permutaties van $X_{m+1}$. Inductie tot $m = 14$:
$G \supseteq A_{15}$ (even permutaties van alle vijftien vakjes),
en $G \subseteq A_{15}$ omdat elke $s_t$ even is: $\langle s_0,
\dots, s_{12}\rangle = A_{15}$.

\textbf{17.} Vragen 13 en 16 geven $A_{15} = \langle s_0, \dots,
s_{12}\rangle \subseteq H$; vraag 9 geeft $H \subseteq A_{15}$.
Dus $H = A_{15}$, van \omterm{def:b2:structures:generated}{orde} $15!/2 = 653\,837\,184\,000$: elke even
herschikking van de vijftien plaatjes is het effect van een
programma.

\textbf{18.} Vraag 8 herleidde $R = \{I = +1\}$ tot het
realiseren van elke even $\rho \in \mathfrak{S}_{15}$ door een
programma: dat is met vraag 17 gedaan. Samen met vraag 6 geeft
dit $\abs R = 16!/2 = 10\,461\,394\,944\,000$. \emph{Twee
klassen:} laat $t_0 = (14\ 15)$ op de \emph{inhouden} werken:
$\varphi(\sigma) = t_0 \circ \sigma$. Een toegestane zet vanuit
$\sigma$ is een toegestane zet vanuit $\varphi(\sigma)$ (het vakje
van de lege plek verandert niet: $(t_0\sigma)^{-1}(16) =
\sigma^{-1}(t_0(16)) = \sigma^{-1}(16)$, en het verschoven vakje
is hetzelfde), en $\varphi(\sigma \circ \tau) = \varphi(\sigma)
\circ \tau$: $\varphi$ beeldt reeksen zetten bijectief op reeksen
zetten af (het is een involutie). Zij klapt $I$ om:
$\varepsilon(t_0\sigma) = -\varepsilon(\sigma)$, bij hetzelfde
vakje voor de lege plek. Bijgevolg beeldt $\varphi$ de klasse $R =
\{I = +1\}$ van $\mathrm{id}$ bijectief af op de klasse van
$\varphi(\mathrm{id}) = \sigma_L$, die dus heel $\{I = -1\}$ is:
precies twee klassen. Dat is de stelling van Johnson en Story.

\textbf{19.} Nummer de vakjes in leesvolgorde en zij $k = 4(i - 1)
+ j$ het vakje van de lege plek. Tel de inversies van $\sigma$
(paren vakjes $x < y$ met $\sigma(x) > \sigma(y)$): paren van twee
plaatjesvakjes leveren $N$; voor de paren waarin de lege plek
voorkomt geldt dat alle vakjes na de lege plek plaatjes $< 16$
dragen en dus geïnverteerd zijn ($16 - k$ paren), terwijl de
vakjes ervóór nooit geïnverteerd zijn. Dus is $\varepsilon(\sigma)
= (-1)^{N + 16 - k} = (-1)^{N + k}$. Omdat $k = 4(i-1) + j \equiv
j \pmod 2$, volgt
\[
I(\sigma) = (-1)^{N + j}\,(-1)^{i + j} = (-1)^{N + i}
= (-1)^{N + r + 1}
\]
met $i = 5 - r$. Volgens vraag 18 is $\sigma$ oplosbaar dan en
slechts dan als $I(\sigma) = +1$, dan en slechts dan als $N + r$
oneven is. Controle: opgelost, $N = 0$, $r = 1$: oneven, dus
oplosbaar; Loyd, $N = 1$, $r = 1$: even, dus onoplosbaar.

\textbf{20.} \emph{Actie:} $e \cdot \sigma = \sigma \circ
\mathrm{id} = \sigma$ en $g \cdot (h \cdot \sigma) = \sigma \circ
h^{-1} \circ g^{-1} = \sigma \circ (gh)^{-1} = (gh) \cdot \sigma$;
bovendien is $\sigma \circ h^{-1}$ opnieuw een stand met de lege
plek thuis ($h$ houdt vakje $16$ vast). \emph{Vrij:} uit $\sigma
\circ h^{-1} = \sigma$ volgt $h^{-1} = \mathrm{id}$ (stel samen
met $\sigma^{-1}$). \emph{Banen = programmaklassen:} vraag 9 zegt
dat de standen die vanuit $\sigma$ met programma's bereikbaar zijn
precies de $\sigma \circ \pi^{-1}$ met $\pi \in H$ zijn: de baan
$H \cdot \sigma$. \emph{Telling:} de vrijheid maakt $h \mapsto h
\cdot \sigma$ injectief, dus telt elke baan $\abs H = 15!/2$
elementen; de $15!$ standen met de lege plek thuis vallen daarmee
uiteen in $15!\,/\,(15!/2) = 2$ banen --- de schaduw, bij lege
plek thuis, van de twee klassen van Johnson en Story.

\textbf{21.} Het rooster van $3 \times 3$ is bipartiet voor de
schaakbordkleuring: elke stap van een wandeling verandert de
kleur, dus heeft elke \emph{gesloten} wandeling een even lengte.
Een gesloten wandeling die elk van de $9$ vakjes precies één keer
bezoekt zou lengte $9$ hebben, oneven: onmogelijk. De constructie
met de grote rondgang uit Deel III is dus niet beschikbaar voor de
achtpuzzel.

\textbf{22.} \emph{Rondgang langs de rand} $9 \to 8 \to 7 \to 4
\to 1 \to 2 \to 3 \to 6 \to 9$ (alle stappen tussen buren; lengte
$8$, even): met de boekhouding van vraag 11 en $c_1 = 8, c_2 = 7,
c_3 = 4, c_4 = 1, c_5 = 2, c_6 = 3, c_7 = 6$ is het effect
\[
\zeta' = (8\ 6\ 3\ 2\ 1\ 4\ 7),
\]
een $7$-cykel die het midden $5$ vasthoudt (de inhoud van $7$ gaat
naar $8$, die van $4$ naar $7$, die van $1$ naar $4$, die van $2$
naar $1$, die van $3$ naar $2$, die van $6$ naar $3$, en die van
$8$ naar $6$). \emph{Hoekrondgang} $9 \to 6 \to 5 \to 8 \to 9$:
effect $(6\ 8\ 5)$ (de inhoud van $5$ gaat naar $6$, die van $8$
naar $5$, en die van $6$ --- geparkeerd in $9$ --- naar $8$). Zet
$y_t = \zeta'^{\,t}(8)$: $y_0 = 8, y_1 = 6, y_2 = 3, y_3 = 2, y_4
= 1, y_5 = 4, y_6 = 7$. Conjugatie (vraag 12) geeft
\[
\zeta'^{\,t}\,(6\ 8\ 5)\,\zeta'^{-t}
= (y_{t+1}\ y_t\ 5) =: T_t \in H_{3\times3},
\]
want $\zeta'$ houdt $5$ vast. De dragers van $T_0 = (y_1\ y_0\ 5)$
en $T_1 = (y_2\ y_1\ 5)$ delen precies $\{y_1, 5\}$: Lemma A geeft
alle even permutaties van $\{y_0, y_1, y_2, 5\}$. Vervolgens voegt
$T_2 = (y_3\ y_2\ 5)$ met Lemma B de letter $y_3$ toe (haar
letters $y_2, 5$ liggen in de huidige verzameling, $k = 4$), en
voegen $T_3, T_4, T_5$ achtereenvolgens $y_4, y_5, y_6$ toe: alle
even permutaties van de acht niet-thuisvakjes liggen in de
programmagroep, die zelf ook uit even permutaties bestaat (het
argument van vraag 9 hangt niet van het bord af). Dus is
$H_{3\times3} = A_8$, en de redenering van de vragen 6, 8 en 18
--- eveneens onafhankelijk van het bord --- laat zien dat de
bereikbare standen precies die met $I = +1$ zijn: de helft van
$9!$, dat wil zeggen $181\,440$.

\textbf{23.} Nummer de vakjes $0, \dots, n-1$ langs de \omterm{def:b2:structures:sn}{cykel}. Een
zet verwisselt de lege plek met een van haar twee buren. Lees de
plaatjes in cyclische volgorde, beginnend net na de lege plek: een
woord $w$ dat de $n - 1$ plaatjes opsomt. De lege plek één stap
vooruit bewegen vervangt $(p, w)$ door $(p + 1, \rho w)$, waarbij
$p$ het vakje van de lege plek is en $\rho$ het woord \omterm{def:b2:structures:generated}{cyclisch} één
plaats roteert; de zet achteruit is de inverse. De \emph{cyclische}
volgorde van de plaatjes (het woord op rotatie na) is dus
invariant. De bereikbaarheidsklasse van $(p, w)$ is de baan van de
afbeelding $g \colon (p, w) \mapsto (p+1, \rho w)$, een element
van \omterm{def:b2:structures:generated}{orde} $\operatorname{lcm}(n, n-1) = n(n-1)$ in het product van
de twee \omterm{def:b2:structures:generated}{cyclische groepen} (translaties van $\Z/n\Z$ en rotaties
van de $n-1$ woordposities), waarbij het kleinste gemene veelvoud
$n(n-1)$ is omdat $\gcd(n, n-1) = 1$: elke klasse telt precies
$n(n-1)$ standen, alle met hetzelfde halssnoer. Aantal klassen:
$n!\,/\,\bigl(n(n-1)\bigr) = (n-2)!$. Voor $n \geq 5$ is $(n-2)! >
2$: de pariteitsinvariant (hoogstens twee klassen) is blind voor
bijna de hele hindernis; de rijkdom van het bord van $4 \times 4$
--- waar de pariteit de \emph{enige} hindernis is --- is een echt
meetkundig feit, geen formeel feit.

\textbf{24.} Beide bakjes hebben de plaatjes in volledig
omgekeerde volgorde, dus is in beide gevallen $N = \binom{15}{2} =
105$ (elk paar plaatjes is geïnverteerd). \emph{Lege plek thuis:}
$r = 1$, $N + r = 106$ even: onoplosbaar. \emph{Lege plek in vakje
$1$:} de lege plek staat in de bovenste rij, $r = 4$, $N + r =
109$ oneven: oplosbaar. Twee bakjes die alleen verschillen in waar
het gat zit, vallen aan weerszijden van de muur.

\textbf{25.} \emph{Morfisme-eigenschap:} zij zet ``één zet = één
\omterm{def:b2:structures:sn}{transpositie}'' om in ``één zet = één tekenwissel'' (vragen 2 en
4), en maakt $I$ zo zet voor zet berekenbaar. \emph{Lagrange:} hij
dwong $6 \mid \abs{\langle s, t\rangle}$ af in Lemma A en bepaalde
de grootte van de nevenklassen bij het uitsluiten van \omterm{def:b2:structures:generated}{orde} $6$
(vraag 14). \emph{Voortbrenging door $3$-cykels:} zij zette ``$H$
bevat genoeg $3$-cykels'' om in ``$H$ bevat heel $A_{15}$'' (vraag
16). \emph{Conjugatie:} zij fabriceerde de vijftien opeenvolgende
$3$-cykels uit één enkele rondgang van $2 \times 2$, meegevoerd
door de grote rondgang (vragen 12--13), en de $3$-cykels $(a\ b\
w)$ in Lemma B. \emph{Metaprincipe:} een invariant bewijst
onmogelijkheid, een expliciete constructie bewijst mogelijkheid,
en een probleem is pas volledig opgelost wanneer de twee grenzen
elkaar raken --- hier op precies de helft.
\end{solution}
